\documentclass[11pt]{article}
\usepackage{varnernotes}

\newcommand{\R}{\mathbb{R}}
\newcommand{\E}{\mathbb{E}}
\newcommand{\Cov}{\operatorname{Cov}}
\newcommand{\Var}{\operatorname{Var}}
\newcommand{\np}{|\mathcal P|}
\newcommand{\bw}{\mathbf w}
\newcommand{\bg}{\mathbf g}
\newcommand{\bal}{\boldsymbol\alpha}
\newcommand{\bbe}{\boldsymbol\beta}
\newcommand{\br}{\mathbf r}
\newcommand{\Tan}{\mathcal T}
\newcommand{\SigSIM}{\hat\Sigma_{g,\mathrm{SIM}}}
\newcommand{\gSIM}{\hat{\bg}_{\mathrm{SIM}}}

\begin{document}

\lecturenotes{6b}{Single Index Model Portfolio Allocation}{Fall 2026}{October 1, 2026}{Jeffrey D. Varner}

\learningobjectives
  {Construct and diagnose SIM portfolio inputs}
  {Assemble the SIM mean vector and covariance matrix from estimated intercepts, betas, residual variances, and training-period market moments; verify that the SIM mean vector reproduces the sample means; identify the residual covariances the SIM omits.}
  {Solve SIM allocation problems}
  {Compute long-only risky-asset and risky/risk-free minimum-variance portfolios from SIM inputs, recover the tangent portfolio from the ray of risky/risk-free solutions, and state where position bounds, a borrowing restriction, or a higher borrowing rate change the picture.}
  {Evaluate what the model changes}
  {Compare SIM and data-driven weights and covariance regret in sample, then evaluate frozen allocations on the 2025 test period without treating one split as a validation of either input set.}

\section{Concept Review: The Single Index Model and Its Covariance}

For asset $i$ in a portfolio $\mathcal P$ of $\np$ risky assets on trading day $t$, with growth rate $g_{i,t}$ (units time$^{-1}$), market growth rate $g_{M,t}$ (SPY; the letter $M$ names the index), intercept $\alpha_i$, market exposure $\beta_i$, and residual $\varepsilon_{i,t}$, the single index model and its population assumptions are
\begin{equation}
  \label{eq:sim}
  g_{i,t}=\alpha_i+\beta_i g_{M,t}+\varepsilon_{i,t},
  \qquad
  \E[\varepsilon_{i,t}]=0,\quad
  \Cov(\varepsilon_{i,t},g_{M,t})=0,\quad
  \Cov(\varepsilon_{i,t},\varepsilon_{j,t})=0\ (i\ne j).
\end{equation}
Stacking the assets, with $\bar g_M=\E[g_M]$, $\sigma_{g,M}^2=\Var(g_M)$, and $D_g=\operatorname{diag}(\sigma_{g,\varepsilon,1}^2,\ldots,\sigma_{g,\varepsilon,\np}^2)$, the mean vector and covariance matrix implied by the model are (L6a)
\begin{equation}
  \label{eq:sim-moments}
  \bar\bg=\bal+\bbe\,\bar g_M,
  \qquad
  \Sigma_g=\sigma_{g,M}^2\,\bbe\bbe^{\mathsf T}+D_g,
\end{equation}
rank one plus diagonal: $2\np+1$ numbers for the covariance and $\np+1$ for the mean in place of $\np(\np+1)/2$ covariances. For fixed weights $\bw$ the linearized one-period portfolio growth rate $g_p=\bw^{\mathsf T}\bg$ has portfolio beta $\beta_p=\bw^{\mathsf T}\bbe$ and variance $\beta_p^2\sigma_{g,M}^2+\bw^{\mathsf T}D_g\bw$. Everything above is a population statement; the inputs below are estimates.

\section{SIM Portfolio Inputs from the Archive}

The L6a archive holds, for each security fitted against SPY on the $N$ trading days of the 2014--2024 training period, the least-squares estimates $\hat\alpha_i$, $\hat\beta_i$, and the residual variance $s^2_{g,\varepsilon,i}=\lVert\br_i\rVert_2^2/(N-2)$, where $\br_i$ is the fitted residual series, together with the sample mean $g'_M$ and sample variance $s^2_{g,M}$ of the market growth rate over the same days. Substituting estimates for population quantities in \eqnref{sim-moments} gives
\begin{equation}
  \label{eq:sim-inputs}
  \gSIM=\hat{\bal}+\hat{\bbe}\,g'_M,
  \qquad
  \SigSIM=s^2_{g,M}\,\hat{\bbe}\hat{\bbe}^{\mathsf T}+\hat D_g,
  \qquad
  \hat D_g=\operatorname{diag}(s^2_{g,\varepsilon,1},\ldots,s^2_{g,\varepsilon,\np}).
\end{equation}
The market moments must be the training-period values; taking them from a later window mixes an in-sample fit with out-of-sample information. All quantities are in growth-rate units (means time$^{-1}$, variances time$^{-2}$, no factor of the observation interval $\Delta t$, one trading day in years here). The covariance of one-period log returns ($\Delta t$ times the growth rates) is $(\Delta t)^2\Sigma_g$, and the GBM covariance rate is $C=\Delta t\,\Sigma_g$; neither belongs in a growth-rate variance formula.

Two exact facts relate these to the data-driven inputs of L5b, the sample-mean vector $\hat{\bg}$ and sample covariance $\hat\Sigma_g$ of the same assets over the same days.

\begin{proposition}{The SIM mean vector is the sample-mean vector}{sim-mean}
A least-squares fit with an intercept has residuals that sum to zero, so $\hat\alpha_i+\hat\beta_i g'_M=g'_i$ for every asset. Hence $\gSIM=\hat{\bg}$ whenever the SIM was fitted on the trading days used for the sample means. The two input sets agree on the reward input.
\end{proposition}

\begin{proposition}{The SIM covariance is the sample covariance with the off-diagonal residual covariances removed}{sim-cov}
Least squares makes each $\br_i$ orthogonal to $\mathbf 1$ and to the market series, so
\begin{equation}
  \label{eq:cov-decomposition}
  \hat\Sigma_{g,ij}=\hat\beta_i\hat\beta_j\,s^2_{g,M}+\widehat{\Cov}(\br_i,\br_j),
  \qquad
  \bigl(\hat\Sigma_g-\SigSIM\bigr)_{ij}=
  \begin{cases}
    \widehat{\Cov}(\br_i,\br_j), & i\ne j,\\[1mm]
    -\lVert\br_i\rVert_2^2/[(N-1)(N-2)], & i=j.
  \end{cases}
\end{equation}
Off the diagonal the difference is exactly the residual covariance; on the diagonal it is the gap between dividing the squared residuals by $N-1$ and by $N-2$, a factor of about $1.0004$ for eleven years of daily data.
\end{proposition}

So off the diagonal the SIM covariance error is the residual covariance, pair by pair (on the diagonal only the negligible degrees-of-freedom factor), and its sign structure is known from L6a. In the examples' universe the residuals of firms in the same industry are positively correlated (banks with banks by several tenths) and those of the technology firms and the banks are negatively correlated by two or three tenths.

\begin{proposition}{The SIM covariance is positive semidefinite}{sim-psd}
For every $\mathbf x\in\R^{\np}$,
\begin{equation}
  \label{eq:sim-psd}
  \mathbf x^{\mathsf T}\SigSIM\mathbf x
  =s^2_{g,M}(\hat{\bbe}^{\mathsf T}\mathbf x)^2+\sum_{i=1}^{\np}s^2_{g,\varepsilon,i}x_i^2\geq0 .
\end{equation}
The matrix is positive definite whenever every residual variance is strictly positive (a sufficient condition; with $s^2_{g,M}>0$ a single zero residual variance is harmless if the corresponding beta is nonzero).
\end{proposition}

Positive definiteness is what the closed forms of L5b require and what makes the variance objective strictly convex. It says nothing about whether the matrix describes future co-movement well.

\section{Risky Asset Portfolios using Single Index Models}

With $g_\star$ the investor's target growth rate, the long-only minimum-variance allocation with SIM inputs is L5b's quadratic program,
\begin{equation}
  \label{eq:risky-qp}
  \begin{aligned}
  \min_{\bw\in\R^{\np}}\quad &\bw^{\mathsf T}\SigSIM\bw
    =\underbrace{s^2_{g,M}(\hat{\bbe}^{\mathsf T}\bw)^2}_{\text{market risk}}
    +\underbrace{\textstyle\sum_i s^2_{g,\varepsilon,i}w_i^2}_{\text{residual risk}}\\
  \text{subject to}\quad
    &\gSIM^{\mathsf T}\bw\geq g_\star,\qquad
    \mathbf 1^{\mathsf T}\bw=1,\qquad
    0\le w_i\le1 .
  \end{aligned}
\end{equation}
The target is a floor: sweeping $g_\star$ upward from the global minimum-variance growth rate traces the GMV portfolio and the efficient branch. Because the two input sets share the mean vector, every difference between the SIM frontier and the data-driven frontier comes from the residual covariances (plus the negligible diagonal factor). For the same weights, by \eqnref{cov-decomposition},
\begin{equation}
  \label{eq:variance-gap}
  \bw^{\mathsf T}\SigSIM\bw
  =\bw^{\mathsf T}\hat\Sigma_g\bw
  -\sum_{i\ne j}w_iw_j\,\widehat{\Cov}(\br_i,\br_j)
  +\sum_i\frac{\lVert\br_i\rVert_2^2}{(N-1)(N-2)}\,w_i^2,
\end{equation}
so the SIM understates the variance of a portfolio concentrated in assets with positively correlated residuals and overstates the variance of one spread across assets with negatively correlated residuals. The two raw frontiers are therefore not directly comparable; the honest comparison evaluates the SIM-optimal weights $\bw_{\mathrm{SIM}}(g_\star)$ under the data covariance and compares with the data-optimal weights $\bw(g_\star)$ at the same target and constraints:
\begin{equation}
  \label{eq:regret}
  \operatorname{regret}(g_\star)
  =\bw_{\mathrm{SIM}}(g_\star)^{\mathsf T}\hat\Sigma_g\,\bw_{\mathrm{SIM}}(g_\star)
  -\bw(g_\star)^{\mathsf T}\hat\Sigma_g\,\bw(g_\star)\geq0 .
\end{equation}
The regret cannot be negative because the data-optimal weights minimize exactly that quantity over the same feasible set. For the thirteen firms of the examples the variance regret runs from about $0.01$ to $1.2$ (time$^{-2}$); the corresponding gap between the two portfolios' growth-rate standard deviations runs from a few thousandths to a tenth of a unit (largest at the concentrated high-growth end), and the two input sets select nearly the same firms with weights that differ by up to seven percentage points.

\begin{remark}{The rank-one inverse}{sim-inverse}
If $\hat D_g$ is positive definite, the Sherman--Morrison identity gives
\begin{equation}
  \label{eq:sim-inverse}
  \SigSIM^{-1}
  =\hat D_g^{-1}
  -\frac{s^2_{g,M}\hat D_g^{-1}\hat{\bbe}\hat{\bbe}^{\mathsf T}\hat D_g^{-1}}
  {1+s^2_{g,M}\hat{\bbe}^{\mathsf T}\hat D_g^{-1}\hat{\bbe}},
\end{equation}
so the closed forms of L5b (GMV and tangent portfolios with short positions allowed) can be evaluated with a diagonal inverse and one inner product rather than a dense matrix inverse. This is why a factor model scales to large universes; the constrained problem \eqnref{risky-qp} still needs the quadratic-program solver.
\end{remark}

\notebooklink
  {L6b: Single Index Model Inputs and Minimum-Variance Portfolios}
  {https://github.com/varnerlab/CHEME-5660-CourseRepository-Fall-2026/blob/main/lectures/week-6/L6b/CHEME-5660-L6b-Example-SIM-MinVar-RA-Fall-2026.ipynb}
  {Assemble the SIM inputs for thirteen firms from the L6a archive, verify Propositions~\ref{prop:sim-mean} and \ref{prop:sim-cov} numerically, trace the long-only frontier from each input set, compare weights and regret, and score the portfolios on 2025 data.}

\section{Risky and Risk-Free Asset Portfolios using Single Index Models}

Let $g_f$ be the growth rate of a horizon-matched risk-free payoff (zero variance over the holding period), $w_f$ its weight, and $\bw$ the risky weights, with $w_f+\mathbf 1^{\mathsf T}\bw=1$. The problem the course package solves is
\begin{equation}
  \label{eq:risky-riskfree-qp}
  \begin{aligned}
  \min_{\bw}\quad &\bw^{\mathsf T}\SigSIM\bw\\
  \text{subject to}\quad
    &\gSIM^{\mathsf T}\bw+(1-\mathbf 1^{\mathsf T}\bw)\,g_f\geq g_\star,\qquad
    0\le w_i\le1 .
  \end{aligned}
\end{equation}
Only the risky weights carry variance and only they are bounded; $w_f=1-\mathbf 1^{\mathsf T}\bw$ is not sign-constrained, so $w_f<0$ (borrowing at $g_f$) is allowed, limited only by the bounds on the risky positions. Eliminating $w_f$ turns the growth constraint into a constraint on the excess growth rate,
\begin{equation}
  \label{eq:excess-return-constraint}
  (\gSIM-g_f\mathbf 1)^{\mathsf T}\bw\geq g_\star-g_f .
\end{equation}

\begin{proposition}{The solutions lie on a ray}{ray}
Let $g_\star>g_f$ be feasible and let $\SigSIM$ be positive definite. Scaling the risky weights by a positive constant scales the excess growth by that constant and the variance by its square, and leaves the constraints $w_i\ge0$ unchanged; hence, as long as no upper bound $w_i\le1$ is active, the growth constraint binds and the minimizers for all targets are one risky direction scaled by the target. Writing $\bw_{\Tan}$ for the point on the ray with $\mathbf 1^{\mathsf T}\bw=1$,
\begin{equation}
  \label{eq:ray}
  \bw(g_\star)=(1-w_f)\,\bw_{\Tan},
  \qquad
  1-w_f=\frac{g_\star-g_f}{\E[g_{\Tan}]-g_f},
\end{equation}
valid until the first bound binds at $(1-w_f)\max_i w_{\Tan,i}=1$, beyond which the boundary kinks. Targets below $\E[g_{\Tan}]$ lend, targets above it borrow, and $\bw_{\Tan}=\bw/\mathbf 1^{\mathsf T}\bw$ for any solution with $\mathbf 1^{\mathsf T}\bw>0$ (at $g_\star=g_f$ the solution is $\bw=\mathbf 0$).
\end{proposition}

\begin{remark}{When a Treasury STRIP is risk-free}{strip-risk}
A zero-coupon Treasury STRIP delivers a specified nominal payment only when its maturity matches the investor's horizon and it is held to maturity; sold earlier, its price carries the interest-rate risk of L2b. ``Risk-free'' here means a model input with a known nominal horizon payoff, not an asset free of every economic risk.
\end{remark}

\subsection*{The capital allocation line and the tangent portfolio (recall from L5b)}

For a risky portfolio $p$ with expected growth $\E[g_p]$ and growth-rate standard deviation $\sigma_{g,p}>0$, the complete portfolio $c$ with fraction $w_f$ risk-free has $\E[g_c]=g_f+(1-w_f)(\E[g_p]-g_f)$ and $\sigma_{g,c}=|1-w_f|\sigma_{g,p}$; eliminating $w_f$ (for $w_f\le1$) gives the capital allocation line
\begin{equation}
  \label{eq:cal}
  \E[g_c]=g_f+\underbrace{\frac{\E[g_p]-g_f}{\sigma_{g,p}}}_{\mathrm{SR}_p}\,\sigma_{g,c},
\end{equation}
whose slope is the Sharpe ratio of $p$. The tangent portfolio $\Tan$ maximizes the Sharpe ratio, and the ray of Proposition~\ref{prop:ray} is its CAL: the excess-growth constraint asks for the least variance per unit of excess growth, so the ray's direction is the long-only maximum-Sharpe portfolio. Under the SIM inputs, with $\hat\beta_{\Tan}=\hat{\bbe}^{\mathsf T}\bw_{\Tan}$,
\begin{equation}
  \label{eq:sim-sharpe}
  \mathrm{SR}_{\Tan}
  =\frac{\hat{\bal}^{\mathsf T}\bw_{\Tan}+\hat\beta_{\Tan}g'_M-g_f}
  {\sqrt{s^2_{g,M}\hat\beta_{\Tan}^2+\bw_{\Tan}^{\mathsf T}\hat D_g\bw_{\Tan}}} .
\end{equation}
When short positions are allowed, $\SigSIM\succ0$, and the GMV portfolio earns more than $g_f$ in expectation, the tangent portfolio has the closed form $\bw_{\Tan}\propto\SigSIM^{-1}(\gSIM-g_f\mathbf 1)$ normalized to sum to one, and \eqnref{sim-inverse} evaluates the inverse. The Sharpe ratio in \eqnref{cal} is the CAL slope in growth-rate units, a one-observation-interval quantity; dividing by the volatility $\sqrt{\Delta t}\,\sigma_{g,p}$ instead ($\Delta t$ the observation interval in years) gives the annualized value $\mathrm{SR}_p/\sqrt{\Delta t}$, larger by $\sqrt{252}\approx15.9$ for daily data. The two rank portfolios identically; from L6b onward the course quotes the annualized value, about $1.3$ for both the SIM and the data-driven tangent portfolios of the examples on the training data.

\begin{figure}[ht]
  \centering
  \begin{tikzpicture}[x=1.05cm,y=0.78cm,>=Latex,font=\sffamily\small]
    \draw[vnink,line width=0.8pt,-{Latex[length=3mm]}] (0,0)--(8.2,0) node[right] {$\sigma_g$};
    \draw[vnink,line width=0.8pt,-{Latex[length=3mm]}] (0,0)--(0,5.8) node[above] {$\E[g]$};
    \draw[vnlink,line width=1.5pt,domain=2.0:6.8,samples=70,smooth]
      plot (\x,{1.65+0.18*(\x-2.0)^2});
    \coordinate (rf) at (0,0.85);
    \coordinate (tan) at (4.65,{1.65+0.18*(4.65-2.0)^2});
    \draw[vncarnelian,line width=1.6pt] (rf)--(tan);
    \draw[vncarnelian,dashed,line width=1.2pt] (tan)--($(rf)!1.55!(tan)$);
    \fill[vnink] (rf) circle (2pt) node[left] {$g_f$};
    \fill[vncarnelian] (tan) circle (2.2pt) node[above left] {$\Tan$};
    \node[text=vncarnelian,rotate=18] at (2.1,2.3) {lending, $0<w_f<1$};
    \node[text=vncarnelian] at (7.0,3.05) {borrowing, $w_f<0$};
    \node[text=vnlink] at (5.0,5.0) {risky frontier};
  \end{tikzpicture}
  \caption{The solid segment is the lending part of the capital allocation line, $0<w_f<1$; the dashed extension borrows at the modeled risk-free rate, $w_f<0$, which the package allows and a no-borrowing constraint removes.}
  \label{fig:constrained-cal}
\end{figure}

\subsection*{Separation, and what breaks it}

Because the tangent portfolio's CAL dominates every other, all mean-variance investors who share the same estimates, the same $g_f$, the same horizon, and the same opportunity set (universe and constraints) in a frictionless market hold the same risky fund $\Tan$ and differ only in $w_f$ (Tobin's two-fund separation, L5b). The result is exact for the unconstrained problem with one lending and borrowing rate, and it survives the long-only constraint on the risky weights alone, which is why the package's ray exists. It breaks under a constraint on the risk-free leg or on dollar positions. No borrowing ($w_f\ge0$) truncates the straight CAL segment at $\Tan$; above it the efficient boundary generally returns to the risky-only frontier, so separation holds only piecewise. Position bounds ($w_i\le1$ here, or a concentration limit $w_i\le u$) kink the ray where they bind; practitioners often impose concentration limits because unconstrained maximum-Sharpe portfolios put extreme weights on a few assets when the mean growth rates are imprecise, and the limits lower the in-sample Sharpe ratio and may improve robustness, a claim to test rather than assume. A borrowing rate above the lending rate kinks the line at $\Tan$.

\notebooklink
  {L6b: SIM Portfolios with a Risk-Free Asset, the Tangent Portfolio, and the Capital Allocation Line}
  {https://github.com/varnerlab/CHEME-5660-CourseRepository-Fall-2026/blob/main/lectures/week-6/L6b/CHEME-5660-L6b-Example-SIM-MinVar-RRFA-Fall-2026.ipynb}
  {Sweep the risky/risk-free problem with SIM inputs, verify Proposition~\ref{prop:ray} numerically, read the tangent portfolio off the ray and check it against the maximum-Sharpe frontier point, compare it with the data-driven tangent portfolio, and score complete portfolios that lend, hold, or borrow on 2025 data.}

\section{What the SIM Inputs Change, and How to Check}

The SIM leaves the reward input unchanged, replaces the sample covariance by a rank-one-plus-diagonal matrix that omits the residual covariances, and does so with $2\np+1$ estimated numbers instead of $\np(\np+1)/2$. For a universe like the examples', where most residual correlations are small (median absolute value $0.14$) even though a few industry pairs reach $0.6$ to $0.7$, that costs little in sample; the cost grows with the residual correlations among the assets the optimizer actually holds. What it buys is not visible in a thirteen-firm example: with hundreds of assets and a few years of daily data the sample covariance has more parameters than the data can pin down and its inverse amplifies the noisiest directions, while the SIM covariance is estimated from a few numbers per asset, is always positive semidefinite, and is interpretable.

Checking follows L5b. Fit on the training window and freeze the universe, the SIM parameters, the market moments, $g_f$ at the decision date, and the weights; evaluate the frozen weights with the buy-and-hold wealth rule on the test year next to simple baselines (equal weights, the GMV portfolio, an index fund); and read the result as one year, one draw, one universe. In the examples the paired SIM and data-driven portfolios earned nearly the same 2025 wealth, which shows that the two input sets selected nearly the same portfolios and cannot rank them, because the gap between them is far smaller than the sampling error in either. The realized risk of the complete portfolios scaled almost exactly with the risky fraction, as the CAL predicts for the linearized model (a buy-and-hold mixture drifts, so not exactly); the realized growth did not match the model, because it depends on whether the tangent portfolio's estimated growth rate persisted.

Finally, the archived parameters carry the standard errors of L6a, and a portfolio inherits that uncertainty nonlinearly: narrow beta intervals do not imply stable weights. Propagating the estimation uncertainty through the covariance into the weights and the portfolio variance is the subject of the optional advanced notebook, and it is the question to ask before any SIM portfolio is traded.

\section*{Optional Advanced Material}

\notebooklink
  {L6b Advanced: Propagating SIM Parameter Uncertainty into a Portfolio}
  {https://github.com/varnerlab/CHEME-5660-CourseRepository-Fall-2026/blob/main/lectures/week-6/L6b/advanced/uncertainty/CHEME-5660-L6b-Advanced-SIM-Portfolio-Uncertainty-Fall-2026.ipynb}
  {Bootstrap the SIM parameters of a six-firm universe two ways, keep each asset's joint draw of intercept, beta, and residual scale, rebuild the SIM covariance for every draw, and read the distributions of a fixed unconstrained global minimum-variance portfolio's growth-rate standard deviation, allocation distance, and variance regret.}

\newpage
\thispagestyle{fancy}
\keytakeaways
  {In this lecture, we assembled the single index model inputs for a portfolio, related them exactly to the data-driven inputs, solved the risky-only and risky/risk-free allocation problems with them, recalled the capital allocation line, the tangent portfolio, and two-fund separation in the form the package computes, and compared the SIM and data-driven portfolios in and out of sample.}
  {The SIM changes only the covariance, and we know how}
  {Fitted with an intercept on the same data, the SIM mean vector equals the sample-mean vector and the SIM covariance equals the sample covariance with the off-diagonal residual covariances removed (plus a negligible degrees-of-freedom factor), so every difference between the two portfolios traces to a residual covariance the model set to zero.}
  {The risky/risk-free problem has one risky answer}
  {With the risk-free fraction free, the growth constraint is a constraint on excess growth and every target selects the same risky direction scaled by $1-w_f$; the fully invested point is the long-only tangent portfolio, the ray is its CAL, and borrowing restrictions, position bounds, or a higher borrowing rate kink or truncate it.}
  {Compare by regret in sample and by frozen weights out of sample}
  {SIM weights evaluated under the data covariance can only lose to the data-optimal weights at the same target, and by little for a universe whose residual correlations are mostly small; one test year checks the weights without ranking the input sets, and parameter uncertainty still has to be propagated into the weights.}
  {Next, in L7b, we let the weights evolve through time and distinguish a rebalancing policy from a one-time allocation.}

\begin{readinglist}
  \item William F. Sharpe, \href{https://doi.org/10.1287/mnsc.9.2.277}{``A Simplified Model for Portfolio Analysis,''} \emph{Management Science} 9(2), 277--293 (1963).
  \item Harry Markowitz, \href{https://www.jstor.org/stable/2975974}{``Portfolio Selection,''} \emph{Journal of Finance} 7(1), 77--91 (1952).
  \item James Tobin, \href{https://doi.org/10.2307/2296205}{``Liquidity Preference as Behavior Towards Risk,''} \emph{Review of Economic Studies} 25(2), 65--86 (1958).
  \item Edwin J. Elton, Martin J. Gruber, Stephen J. Brown, and William N. Goetzmann, \emph{Modern Portfolio Theory and Investment Analysis}, 9th ed. (2014), chapters on index models and the capital allocation line.
\end{readinglist}

\vfill
{\sffamily\footnotesize\color{vnmuted}\textbf{Educational use and risk notice.} These notes are for instruction only and do not constitute investment advice. Financial decisions involve risk, and model outputs depend on assumptions.}

\end{document}
